\chapter{Atomic data}
\label{ch:data}

By default the module computes the atomic physics of every element at start-up from a screened-hydrogenic shell
model, for the photon energy of the pulse (\opt{mcmd-atomic-data = hydrogenic}, Section~\ref{sec:builtin}); no input
file is needed.
With \opt{mcmd-atomic-data = files} it reads the same quantities from plain-text tables in \file{./Atomic\_data/} of
the run directory instead.
This chapter specifies the tables; \file{ATOMIC\_DATA.md} in the repository is the same specification in short form,
and \file{tools/check\_atomic\_data.py} checks a directory against it.

\section{The built-in data}
\label{sec:builtin}

The built-in data are a screened-hydrogenic model of shells only~\cite{scott2010}, with the screening constants of
Marchand, Caille and Lee~\cite{marchand1990} and the level energies scaled to the NIST ionization
energies~\cite{nist}:
\begin{itemize}
  \item states: every \code{K L M} of every charge state ($K \le 2$, $L \le 8$, $M \le 18$), so elements up to
    $Z = 18$ (H to Ar except He); Ca, Fe, Ni and I need \opt{mcmd-atomic-data = files};
  \item energies: the screened-hydrogenic binding energy of each configuration (with the leading relativistic
    terms), the excitation energy scaled by the ratio of the NIST to the hydrogenic ionization energy of the charge
    state its lowest inner-shell hole belongs to;
  \item photoionization out of each occupied shell: Kramers' cross section with a constant Gaunt factor,
    $\sigma \propto q^4 n^{-5} E^{-3}$ per electron ($q$ the screened charge of the shell);
  \item fluorescence: the Einstein $A$ of each radiative transition into a lower shell, from hydrogenic oscillator
    strengths scaled with the screened charges;
  \item Auger: the hydrogenic formula of Scott and Hansen~\cite{scott2010} (Gaunt factor 0.12), summed per final state;
  \item statistical weights $g = \binom{2}{K}\binom{8}{L}\binom{18}{M}$.
\end{itemize}
It takes well under a second.
The model is the one of the CRETIN screened-hydrogenic models the published results were computed with; it reproduces
their rates to $\sim$10\,\% and, on lysozyme at 1000\,eV, their explosion: charge to 0.7\,\%, escape to 1.3\,\%,
directions within the seed noise, ion kinetic energy 3.5\,\% lower (Section~\ref{sec:accuracy}).
There are no collisional fit data: electron-impact ionization comes from Lotz's formula (\opt{mcmd-coll-data = lotz}),
and collisional excitation (\opt{mcmd-coll-excitation}) and \opt{mcmd-coll-data = files} need
\opt{mcmd-atomic-data = files}; \code{grompp} stops otherwise.
The run writes the data it used to \file{simulation\_output/Atomic\_data/}, in the format below; copied to
\file{./Atomic\_data} and run with \opt{mcmd-atomic-data = files} they give the identical run.

\section{Which files}

A file is \file{<table>\_<SYMBOL>.txt}, with the symbol in upper case as in Table~\ref{tab:elements}
(\file{energy\_levels\_C.txt}, \file{rate\_transitions\_CL.txt}).
Only the elements present in the system are needed.

\begin{table}[H]
\centering\small
\begin{tabularx}{\textwidth}{@{}>{\raggedright\arraybackslash}p{0.4\textwidth} X X@{}}
\toprule
file & needed when (\opt{mcmd-atomic-data = files}) & what for \\
\midrule
\file{energy\_levels\_X.txt} & always & outermost binding energy per state \\
\file{rate\_transitions\_X.txt} & always & photoionization, Auger, fluorescence \\
\file{total\_energies\_X.txt} & charge transfer or collisions on (both default) & electron energies, charge-transfer
  energies, collisional transition energies \\
\file{statistical\_weight\_X.txt} & collisions on & detailed balance \\
\file{collisional\_ionization\_table\_X.txt} or \file{collisional\_parameters\_X.txt} &
  only with \opt{mcmd-coll-data = files} & electron-impact ionization \\
\file{collisional\_excitation\_table\_X.txt} or \file{collisional\_excitation\_X.txt} &
  only with \opt{mcmd-coll-excitation = 1} & collisional (de-)excitation \\
\bottomrule
\end{tabularx}
\end{table}
Only \file{rate\_transitions\_X.txt} depends on the photon energy (the photoionization cross sections): regenerate it
for every \opt{mcmd-pulse-photon-energy}.

\paragraph{General rules.}
One record per line, whitespace separated; \code{;} separates the initial state from the final states in the
transition tables; no comments or header (except the tables' temperature line).
Lines must be shorter than 1024 characters.
Energies in eV, rates in s$^{-1}$, cross sections in cm$^2$.
Where a state appears twice, the first line counts.

\section{The tables}

\paragraph{\file{energy\_levels\_X.txt}:} \code{K L M E}, the binding energy of the state's \emph{outermost}
electron.
Autoionizing states (excitation above the ionization potential of their charge state) get \code{1e14}.
Used as the charge-transfer donor energy and the field-ionization threshold.

\paragraph{\file{total\_energies\_X.txt}:} \code{K L M T}, the energy to remove all of the state's electrons,
$T = \sum(\text{ionization potentials up to this charge state}) - \text{excitation energy}$; defined for autoionizing
states too, and \code{0 0 0 0.0}.
It must reproduce \file{energy\_levels} for every bound state, $E(c) = T(c) - T(c\text{ minus its outermost
electron})$, and cover every reachable state.
It gives the kinetic energy of emitted electrons (photoelectron $E_{\mathrm{ph}} - (T_{\mathrm{from}}-T_{\mathrm{to}})$,
Auger $T_{\mathrm{to}} - T_{\mathrm{from}}$, both must be positive), the binding energy of any shell for charge
transfer, and collisional transition energies.

\paragraph{\file{rate\_transitions\_X.txt}:} one line per initial state, then one \code{;}-terminated token per final
state with five fields:
\begin{lstlisting}
K L M ;K' L' M' rate type ;K'' L'' M'' rate type ;...
1 1 0 ;1 0 0 7.416614e-21 2 ;0 1 0 0.0 2 ;2 0 0 7.291753e+11 1 ;      (oxygen, 600 eV)
\end{lstlisting}
Type 0: Auger--Meitner decay, rate in s$^{-1}$, one electron fewer.
Type 1: fluorescence, rate in s$^{-1}$, same electron count.
Type 2: photoionization, cross section in cm$^2$, one electron fewer.
The type column is mandatory: many Auger and photoionization transitions have the same occupation change.
Every state reachable through any transition needs a line of its own, also without transitions (\code{K L M ;});
only \code{0 0 0} may be absent.
A zero rate is inactive; two entries to the same final state are two channels.

\paragraph{\file{statistical\_weight\_X.txt}:} \code{K L M g}, the degeneracy $g = \binom{2}{K}\binom{8}{L}\binom{18}{M}$.

\paragraph{Collisional data} (only for \opt{mcmd-coll-data = files} and collisional excitation):
\begin{itemize}
  \item \file{collisional\_ionization\_table\_X.txt}, \file{collisional\_excitation\_table\_X.txt}: a first line
    \code{T\_e t1 ... tn} (eV, increasing, full precision), then \code{K L M K' L' M' k1 ... kn}: the full Maxwellian
    rate coefficient $\langle\sigma v\rangle$ (cm$^3$\,s$^{-1}$, Boltzmann factor included) at each temperature.
    Any exponent is allowed (\code{3.1e-2650}); a \code{0} means too small to write and is extrapolated.
    The tables are generator independent: FLYCHK, FAC, BEB or Lotz can supply them.
  \item \file{collisional\_parameters\_X.txt} (\code{K L M ; K' L' M' c0 c1 c2 c3 s0 ; ...}) and
    \file{collisional\_excitation\_X.txt} (\code{... c0 c1 c2 c3 a0 ...}): the CRETIN Sampson and \code{colex2} fits,
    used where no table is given (formulas in Chapter~\ref{ch:coll}).
\end{itemize}

\section{Generating the data}
\label{sec:generators}

Two generators write all the tables above, for \opt{mcmd-atomic-data = files}.

\paragraph{From XATOM} (\file{tools/make\_atomic\_data\_xatom.py OUTDIR ENERGY [--elements H,C,N,O,S] [--validate]}).
XATOM~\cite{jurek2016,son2011} is an open-source (GPLv3) Hartree--Fock--Slater atomic code; its output is not covered
by its license, so the generated data can be distributed freely.
The script runs \code{xatom} once per configuration ($\sim$20\,s for H, C, N, O, S on 12 cores):
\begin{itemize}
  \item states: every \code{K L M} up to the ground-state occupation --- photoionization, Auger, fluorescence and charge
    transfer never leave that set;
  \item each state is computed as its aufbau subshell configuration (2s before 2p, 3s before 3p), and XATOM's subshell
    final states are summed back onto \code{K L M};
  \item energies from the orbital eigenvalues,
    $T(c) = T(c\text{ minus its outermost electron}) - \varepsilon_{\mathrm{outermost}}(c)$: XATOM's total energies are
    tens of eV off (N$^{5+}$\,1s$^2$: 1267\,eV against 1218), its eigenvalues are what it prices transitions with;
  \item core and valence photoionization (\code{P}, \code{V}) become type 2, \code{A} type 0, \code{F} type 1;
    photoexcitation is dropped; channels these energies close and intra-shell ones are dropped and counted;
  \item elements up to $Z = 18$ (H to Ar except He); Ca, Fe, Ni are refused (their 4s electrons are
    counted in M).
\end{itemize}

\paragraph{From CRETIN models} (\file{tools/make\_atomic\_data\_cretin.py MODELS OUTDIR ENERGY [--tables] [--validate]}).
Reads CRETIN \file{.z01} atomic models and writes all tables including the Sampson and \code{colex2} fits;
\code{--tables} also writes them as temperature tables (\file{tools/make\_collisional\_tables.py}).
CRETIN is not open source.

\paragraph{Checking.}
\code{tools/check\_atomic\_data.py Atomic\_data --photon-energy E} reports everything the reader would stop on or
misread as an error and everything legal but suspicious (rates out of range, closed channels, unreachable states) as a
warning; it exits 1 on errors.

\section{Accuracy of the data}
\label{sec:accuracy}

Against Scofield's subshell photoionization cross sections~\cite{scofield1973} and Krause's K-shell widths and
fluorescence yields~\cite{krause1979} (via xraylib):
\begin{table}[H]
\centering\small
\begin{tabular}{lccc}
\toprule
ground state, 600\,eV, kb & Scofield & XATOM & CRETIN \\
\midrule
C K / L & 167.8 / 7.86 & 167.8 / 7.86 & 179.5 / 1.61 \\
N K / L & 275.9 / 14.4 & 276.0 / 14.4 & 317.6 / 4.13 \\
O K / L & 411.3 / 24.5 & 411.5 / 24.5 & 507.9 / 9.05 \\
S L / M & 445.9 / 26.6 & 444.7 / 26.5 & 419.2 / 2.19 \\
\bottomrule
\end{tabular}
\end{table}
XATOM agrees with Scofield to 0.3\,\% and with Krause for the S K hole (width 0.56 against 0.52\,eV, fluorescence yield
0.089 against 0.080).
The CRETIN models' K-shell cross sections are 7--23\,\% high, their valence ones 3--12$\times$ low, and their S K-hole
Auger rate 3.4$\times$ low.
XATOM's weak point is the outermost binding energy of neutral C and N ($\sim$20\,\% low, open 2p shells).
On lysozyme at 600\,eV, XATOM data give a 15\,\% higher ion kinetic energy than CRETIN data, 1\,\% more charge, and the
same ion directions; at 1000\,eV (the example) 39\,\% higher, with half the Auger decays and twice the valence
photoionizations.

The built-in data are of the CRETIN kind: core photoionization within 10\,\% of the CRETIN models, valence within
5\,\%, Auger and fluorescence per configuration within $\sim$10\,\%, so with the same deviations from the reference
data; their K-hole Auger rates are 1.5--2.5$\times$ slower than CRETIN's, and the constant Gaunt factor makes the
photoionization rise above CRETIN's towards 10\,keV (C, N, O core 1.3$\times$ at 12\,keV).
